\documentclass[11pt]{article}
\usepackage[margin=1in]{geometry}
\usepackage{amsmath,amssymb,amsthm,hyperref}
\newtheorem{proposition}{Proposition}
\newtheorem{definition}{Definition}
\title{Labor Heterogeneity and Reallocation: Equilibrium and Joint AI Supply}
\author{Working derivations for Oliver Pardo}
\date{22 November 2026}
\begin{document}
\maketitle
I extend the original equilibrium, labelled \texttt{def:rewrite-equilibrium}
in the manuscript. This working note is not yet inserted into the paper.

\section{Environment and equilibrium}
Two competitive activities produce the same final good,
\begin{equation}
 Y=Y_S+Y_C,\quad Y_j=K_j^\alpha Z_j^{1-\alpha},\quad
 Z_j=[\omega_{Lj}(AL_j)^{d_j}+\omega_{Xj}X_j^{d_j}]^{1/d_j},\quad
 d_j=(\sigma_j-1)/\sigma_j,
 \label{eq:extension-production}
\end{equation}
with $0<\sigma_C<1<\sigma_S$. Both products have price one.
Positive CES weights sum to one within each activity. The capital share is
the same in both activities. Capital can move freely, with $K_S+K_C=K$.
In the sector-specific benchmark, $L_j=\ell_jN$, with fixed positive shares
summing to one. In the mobile benchmark, identical workers choose their
activity, and $L_S+L_C=N$. These are alternative labor-market closures.

The household pools labor income and owns capital and the developer. It has
the original preferences $\int_0^\infty e^{-\rho t}N\log(C/N)\,dt$,
population growth $n$, and labor-augmenting growth $\gamma$. Its budget is
\[
 \dot K=rK+w_SL_S+w_CL_C+\Pi-C.
\]
The Euler equation remains $\dot C/C=n+r-\rho$, and the household
transversality condition is unchanged. I do not introduce asset-ownership
heterogeneity or separate consumption decisions by worker type.

A single integrated developer supplies identical AI services at a common
linear price $p_X$. It purchases $U=(X_S+X_C)/B$ units of compute and
research compute $M$. Efficiency follows the original law
\[
 \dot B=\chi(BM)^\eta(1-B/\overline B).
\]
The developer takes aggregate $K,A,N,r$ as given, as in the original model.
Its within-date demand incorporates allocation of capital and AI services
between activities. The sector-specific labor supplies are fixed in that
demand. Under mobile labor, labor reallocation must also be incorporated.
There is no discriminatory pricing or access fee in this benchmark.

\begin{definition}[Extended equilibrium]
An equilibrium is a trajectory of finite allocations and prices satisfying
the original initial conditions and exogenous laws, household optimality,
profit maximization by both final-good activities, and global optimality
of the developer's discounted net profit. Capital and AI markets clear,
the selected labor-market closure holds, and
\begin{equation}
 C+\dot K+\delta K+(X_S+X_C)/B+M=Y.
 \label{eq:extension-resources}
\end{equation}
Under mobile labor, an employed activity pays the maximum available wage.
The developer maximizes operating profit at each date and then solves the
original research problem with that optimized profit function.
Both transversality requirements are retained. First-order conditions alone
do not define global developer optimality.
\end{definition}

Euler's identity in each active activity gives
$Y_j=(r+\delta)K_j+w_jL_j+p_XX_j$.
Summing and using $\Pi=p_X(X_S+X_C)-(X_S+X_C)/B-M$ establishes
\eqref{eq:extension-resources} from the household budget. No AI payment
is counted as an additional resource cost beyond compute and research.

\section{Reduction with sector-specific labor}
Write $H_j=AL_j$ and $X=X_S+X_C$.
At given positive composites, maximizing the sum of outputs over capital
gives $K_j/K=Z_j/(Z_S+Z_C)$, since the capital marginal products must
coincide. This also gives the correct boundary allocation if a composite
is zero. Substitution gives $Y=K^\alpha(Z_S+Z_C)^{1-\alpha}$.
For a given total supply $X$, define
\begin{equation}
 G(X)=\max_{x_S+x_C=X,\ x_j\geq0}
       \{Z_S(H_S,x_S)+Z_C(H_C,x_C)\}.
 \label{eq:extension-G}
\end{equation}
Each CES is concave and strictly concave in its AI argument for positive
labor and positive AI. Thus the allocation is unique. If both activities
use AI, their marginal composite products equal a common $\mu=G'(X)$.
At a zero quantity, use the corresponding right-derivative inequality.

Let $c_S=\omega_{XS}^{\sigma_S/(\sigma_S-1)}<1$ and
$b_C=\omega_{XC}^{\sigma_C/(\sigma_C-1)}>1$.
The allocation can be parameterized explicitly by $\mu\in(c_S,\infty)$:
\begin{equation}
 x_j(\mu)=H_j
 \left[\frac{\omega_{Lj}}
 {(\mu/\omega_{Xj})^{\sigma_j-1}-\omega_{Xj}}\right]^{\sigma_j/(\sigma_j-1)}.
 \label{eq:extension-sector-demand}
\end{equation}
For $S$ this expression applies throughout that interval.
For $C$ use it only at $\mu<b_C$, and set $x_C=0$ at $\mu\geq b_C$.
The formula follows by solving $Z_{j,X}=\mu$ and substituting into the
CES identity. Total $X(\mu)=x_S(\mu)+x_C(\mu)$ is continuous and
strictly decreasing from infinity to zero. Hence every positive supply
has a unique allocation. The complementary activity can be closed at low
AI supplies. Its right marginal product at zero AI is finite, unlike that
of the substituting activity.

The common marginal value of composites and inverse AI demand are
\begin{equation}
 Q=(1-\alpha)(K/G)^\alpha,\qquad P(X)=QG'(X)=Q\mu.
 \label{eq:extension-inverse}
\end{equation}
Demand is strictly decreasing. Away from entry kinks,
\begin{equation}
 e_X=\alpha\frac{XG'}G-\frac{XG''}{G'},\qquad
 G''=\left(\sum_{j:x_j>0}\frac1{Z_{j,XX}}\right)^{-1}.
 \label{eq:extension-elasticity}
\end{equation}
To derive the second equation, differentiate $Z_{j,X}(x_j)=\mu$ and
sum $dx_j/d\mu$. Log-differentiating \eqref{eq:extension-inverse} gives
the first. At a kink use one-sided derivatives, not this smooth formula.

\section{Global within-date monopoly choice}
\begin{proposition}[Global service supply]
At fixed positive $K,H_S,H_C,B$, operating profit
\begin{equation}
 \pi(B)=\max_{X\geq0}\{P(X)X-X/B\}
 \label{eq:extension-monopoly}
\end{equation}
has a positive finite global maximizer. Its set of maximizers is compact.
Every smooth maximizer satisfies $P(X)(1-e_X)=1/B$.
The global solution is obtained by comparing objective values over the
scalar parameterization \eqref{eq:extension-sector-demand}, including
the entry point $\mu=b_C$ and any smooth stationary candidates.
Neither uniqueness nor sufficiency of the first-order condition is claimed.
\end{proposition}
\begin{proof}
As $X\downarrow0$, only $S$ uses AI. Its composite has a positive limit
and marginal product proportional to $X^{-1/\sigma_S}$. Thus $G$ has
a positive limit and revenue is asymptotically proportional to
$X^{1-1/\sigma_S}$. Revenue extends continuously to zero at $X=0$ and
dominates the linear compute cost at sufficiently small positive $X$.

As $X\to\infty$, $\mu\downarrow c_S$, $x_C$ converges to a finite
positive value, $x_S/X\to1$, $G/X\to c_S$, and $G'\to c_S$.
Revenue is therefore of order $X^{1-\alpha}$, while cost is linear.
Profit tends to minus infinity. Continuity gives a positive finite global
maximum and a compact maximizer set. Differentiation at a smooth maximizer
gives the stated necessary condition. The scalar parameterization is
one-to-one, so maximizing over that parameter is globally equivalent.
\end{proof}

Under competitive provision, solve $P(X)=1/B$. Strictly decreasing demand
with endpoint limits infinity and zero gives one positive solution.
Every monopoly maximizer has positive profit, hence $P(X)>1/B$ and less
total AI supply than competition at the same $K,H_S,H_C,B$.
Each sector's quantity is nondecreasing in total supply, by the $\mu$
parameterization. The monopoly therefore weakly reduces each activity's
AI use, strictly in each activity that is active at both allocations.

\subsection{Why the original uniqueness proof cannot simply be reused}
For $\sigma_C<1/2$, the complementary CES has $Z_{C,XX}(0+)=0$.
When it begins using AI, $G''$ changes from $Z_{S,XX}<0$ to a right
limit of zero. The marginal revenue $P+XP'$ therefore jumps upward.
If compute cost lies between the two one-sided marginal revenues,
profit can have a local maximum on each side of entry.

A static numerical example uses $\alpha=0.33$, $\sigma_S=1.5$,
$\sigma_C=0.25$, $K=H_S=H_C=1$, all CES weights $1/2$, and
$B=1.2990442402452336$. At entry the one-sided marginal revenues are
approximately $0.50655$ and $1.03305$, while compute cost is $0.76980$.
There are two smooth local maxima, approximately
\[
 (X,\pi)=(0.04325,0.04583),\qquad (0.27516,0.05924).
\]
The larger-supply local maximum has higher profit. These computations
disprove sufficiency of an arbitrary first-order solution, not uniqueness
of the global optimum. The accompanying Python check records the example
and verifies both price derivatives and capital aggregation.

At a differentiability point of $\pi(B)$, the original envelope remains
$\pi_B=X^*/B^2$. If maximizers differ, differentiability must be checked
before using the original smooth research costate equation.
The dynamic verification lemma and the equilibrium existence proof are
not established by the static result above.

\section{Next proof obligations}
I must establish the global selection relevant to long-run equilibria,
then developer research optimality and existence near the candidate limits.
Under mobile labor, replace \eqref{eq:extension-G} by maximization over
both AI and labor allocation. That separate problem must be solved before
stating a mobile-labor equilibrium proposition. The original manuscript,
its simulations and its closed sections remain unchanged.
\end{document}
